% sections/ground_truth.tex — The Ground Truth Problem
\section{The Ground Truth Problem in Enterprise Knowledge Systems}
\label{sec:ground_truth}

This section formalizes the fundamental challenge that motivates
\system{}: the inherent limitations of recovering complete ground truth from
observed enterprise data alone, and how this motivates the forward-generation
approach.

% ─────────────────────────────────────────────────────────────────────
\subsection{The Inverse Problem of Knowledge Recovery}
\label{sec:gt_inverse}

Enterprise knowledge systems for audit analytics are typically built by
a process we term \emph{inverse knowledge recovery}: given observed data
$D$ (journal entries, document flows, master records), reconstruct the
true generative process $\mathcal{G}$ that produced it.  We formalize
this problem and show it is fundamentally ill-posed.

\begin{definition}[Enterprise Data Observation]
\label{def:observation}
An \emph{enterprise data observation} is a tuple $D = (J, M, R, T)$
where:
\begin{itemize}
  \item $J = \{j_1, \ldots, j_N\}$ is a set of journal entries, each
    with line items $j_k = \{l_{k,1}, \ldots, l_{k,n_k}\}$ where each
    line item specifies an account, amount, and debit/credit indicator;
  \item $M$ is a set of master data records (vendors, customers,
    materials, employees);
  \item $R$ is a set of document references linking records across tables;
  \item $T$ is a set of timestamps and period assignments.
\end{itemize}
\end{definition}

\begin{definition}[Accounting Network]
\label{def:accounting_network}
Given a journal entry $j$ with $n$ credit line items and $m$ debit line
items, an \emph{accounting network} maps each credit-to-debit flow as a
directed edge in a graph $G = (V, E)$ where nodes $V$ are GL accounts
and edges $E$ carry monetary amounts.  The edge construction requires
a \emph{matching} of credit amounts to debit amounts such that flows are
conserved.
\end{definition}

The central difficulty is that the mapping from journal entries to
accounting network edges is not unique.  For a single journal entry
with $n$ credit and $m$ debit line items, the number of valid
matchings is given by the surjective function count.

\begin{proposition}[Combinatorial Infeasibility of Ground Truth Recovery]
\label{thm:combinatorial}
Let $j$ be a journal entry with $n$ credit line items and $m$ debit
line items.  The number of distinct accounting network edge
configurations consistent with $j$ is bounded below by:
\begin{equation}
\mathcal{A}(n, m) = \sum_{k=0}^{m} (-1)^{m-k} \binom{m}{k} k^n
\label{eq:stirling_matching}
\end{equation}
which corresponds to the number of surjections from the credit set to
the debit set~\citep{ivertowski2024hardware}.  For an enterprise with
$N$ journal entries, the total configuration space is:
\begin{equation}
|\mathcal{S}| = \prod_{j=1}^{N} \mathcal{A}(n_j, m_j)
\label{eq:total_configs}
\end{equation}
\end{proposition}

\begin{proof}
The edge construction problem for a single entry requires assigning
each of $n$ credit line items to at least one of $m$ debit line items
such that total flows balance.  This is a surjection from the set of
credits to the set of debits.  The number of surjections from an
$n$-element set to an $m$-element set is $m! \cdot S(n,m)$ where
$S(n,m)$ is the Stirling number of the second
kind~\citep{ivertowski2024hardware}:
\[
S(n, m) = \frac{1}{m!} \sum_{k=0}^{m} (-1)^{m-k} \binom{m}{k} k^n.
\]
Multiplying through gives
$\mathcal{A}(n,m) = \sum_{k=0}^{m}(-1)^{m-k}\binom{m}{k}k^n$,
which is \Cref{eq:stirling_matching}.
Since journal entries are independent, the total configuration space is
the product over all entries.

For concrete scale: in the 155 real-world datasets analyzed
by~\citet{ivertowski2024hardware}, only 60.7\% of entries have $n = m = 1$
(trivial 1:1 matching).  The remaining 39.3\% require non-trivial matching.
For a modest entry with $n = 4$ credits and $m = 3$ debits,
$\mathcal{A}(4,3) = 36$.  An enterprise with $10^5$ such non-trivial
entries has $|\mathcal{S}| > 36^{10^5} \approx 10^{155{,}630}$
possible configurations---a number that dwarfs the number of atoms in
the observable universe ($\sim 10^{80}$).
\end{proof}

\begin{remark}
\Cref{thm:combinatorial} counts discrete credit-to-debit
\emph{assignments}.  When monetary amounts must also be split across
edges, the configuration space becomes continuous and uncountably
infinite, making the discrete count a lower bound on the true
ambiguity.
\end{remark}

\begin{corollary}
\label{cor:verification}
For any enterprise dataset with a non-trivial fraction of multi-line
journal entries, enumerating all consistent accounting network
configurations is computationally infeasible.  Consequently, without
additional information beyond the journal entries themselves,
distinguishing the true configuration from the exponentially many
alternatives is not possible by exhaustive search.
\end{corollary}

% ─────────────────────────────────────────────────────────────────────
\subsection{Systematic Error Propagation}
\label{sec:gt_errors}

The combinatorial infeasibility of \Cref{thm:combinatorial} means that
knowledge recovery must rely on heuristics and assumptions.  We now show
that when the source data contains systematic errors, these propagate
through any reconstruction pipeline with high probability.

\begin{definition}[Systematic Error]
\label{def:systematic_error}
A \emph{systematic error} in enterprise data is a consistent
deviation from the true state that affects multiple records in the same
direction or pattern.  Unlike random errors, systematic errors do not
cancel under aggregation and may be invisible to internal consistency
checks.
\end{definition}

\begin{observation}[Internal Consistency Does Not Imply Correctness]
\label{obs:consistency}
A dataset $D$ may satisfy all standard accounting invariants---balanced
entries (\Cref{eq:balanced}), the accounting equation ($A = L + E$),
subledger reconciliation, document chain integrity---while containing
systematic errors in amounts, entity classifications, period
assignments, or relationship structures.  Internal consistency is
necessary but not sufficient for correctness.
\end{observation}

\begin{proposition}[Error Persistence in Multi-Stage Processes]
\label{prop:error_persistence}
Consider an enterprise process with $k$ sequential stages
(e.g., PO $\to$ GR $\to$ Invoice $\to$ Payment $\to$ GL Posting
$\to$ Financial Statement, giving $k = 6$).  Let each stage have
an independent probability $\varepsilon_d$ of detecting and correcting a
systematic error introduced at an earlier stage.  The probability that
the error reaches the final stage undetected is:
\begin{equation}
\Prob[\text{undetected}] = (1 - \varepsilon_d)^{k-1}
\label{eq:error_persistence}
\end{equation}
For typical detection rates $\varepsilon_d \in [0.01, 0.05]$ and $k = 6$:
\begin{center}
\begin{tabular}{ccc}
\toprule
$\varepsilon_d$ & $k$ & $\Prob[\text{undetected}]$ \\
\midrule
0.01 & 6 & 95.1\% \\
0.02 & 6 & 90.4\% \\
0.05 & 6 & 77.4\% \\
\bottomrule
\end{tabular}
\end{center}
Even with a 5\% per-stage detection rate, over three-quarters of
systematic errors survive the full processing chain.
\end{proposition}

\begin{proof}
Each stage provides an independent opportunity to detect the error.
The error survives stage $i$ with probability $(1 - \varepsilon_d)$.
Since stages are independent, the probability of surviving all $k-1$
stages after introduction is $(1 - \varepsilon_d)^{k-1}$.
\end{proof}

\begin{remark}
The independence assumption in \Cref{prop:error_persistence} is
optimistic: in practice, detection capabilities are often correlated
across stages (e.g., a subtle classification error that evades one
stage is likely to evade similar checks at subsequent stages).  The
actual persistence probability for systematic errors may therefore
be higher than $(1 - \varepsilon_d)^{k-1}$, making this a conservative
estimate of the problem's severity.
\end{remark}

% ─────────────────────────────────────────────────────────────────────
\subsection{Inter-Dimensional Complexity}
\label{sec:gt_interdimensional}

The difficulty of the ground truth problem is compounded by the
\emph{inter-dimensional} nature of enterprise data errors.  We identify
five dimensions along which errors manifest, and show that their
interactions create a complexity class that exceeds any single-dimension
analysis.

\begin{definition}[Knowledge Dimensions]
\label{def:dimensions}
Enterprise audit knowledge spans five orthogonal dimensions:
\begin{enumerate}
  \item \textbf{Topological} ($\mathcal{D}_T$): Account classifications,
    entity relationships, document reference chains, graph connectivity.
  \item \textbf{Statistical} ($\mathcal{D}_S$): Amount distributions,
    Benford's-law compliance, correlation structures, outlier patterns.
  \item \textbf{Temporal} ($\mathcal{D}_\tau$): Posting dates, period
    assignments, processing lags, seasonal patterns, regime changes.
  \item \textbf{Normative} ($\mathcal{D}_N$): Accounting standards
    compliance, audit requirements, internal control effectiveness,
    regulatory adherence.
  \item \textbf{Organizational} ($\mathcal{D}_O$): Approval hierarchies,
    segregation of duties, departmental structure, intercompany
    relationships.
\end{enumerate}
\end{definition}

\begin{observation}[Cross-Dimensional Error Amplification]
\label{thm:amplification}
Let $e$ be a systematic error in dimension $\mathcal{D}_i$ that affects
$n_i$ records.  If the error induces secondary effects in $d$ additional
dimensions, each affecting a multiplicative factor $\alpha_j \geq 1$
records, then the total error footprint is:
\begin{equation}
|e^*| = n_i \cdot \prod_{j=1}^{d} \alpha_j
\label{eq:amplification}
\end{equation}
For typical enterprise processes where $d \geq 2$ and $\alpha_j \geq 2$,
a single-source error affecting 100 records can propagate to $\geq 400$
affected data points across dimensions.
\end{observation}

\paragraph{Illustrative example.}
Consider a vendor master data error (topological dimension) that
assigns incorrect GL account codes to a vendor.  This produces:
(a)~statistical anomalies in the affected accounts' distributions
($\alpha_S \geq 1$), (b)~temporal artifacts from misclassified period
assignments ($\alpha_\tau \geq 1$), and (c)~normative violations in
subledger reconciliation ($\alpha_N \geq 1$).  Each propagation
multiplies the affected record count.  The product follows from the
independence of propagation paths across dimensions.
The multiplicative model provides a useful upper-bound heuristic
for estimating cross-dimensional impact, though actual propagation
depends on the specific error type and process dependencies.

\begin{remark}
The inter-dimensional nature of errors explains why single-dimension
audit analytics---Benford tests alone, or control testing alone, or
process mining alone---systematically underperform.  Effective audit
analytics must operate across all five dimensions simultaneously,
which requires training data that embeds knowledge across all five
dimensions.  This is precisely what \system{}'s reference knowledge
graphs provide.
\end{remark}

% ─────────────────────────────────────────────────────────────────────
\subsection{Three Layers of Audit Knowledge}
\label{sec:gt_layers}

Having established that inverse recovery is infeasible and error-prone,
we now define the \emph{forward generation} approach through three
layers of audit knowledge.  \system{} generates data that explicitly
encodes all three layers, producing reference knowledge graphs with
full provenance.

\begin{definition}[Three-Layer Knowledge Model]
\label{def:three_layers}
A \emph{reference knowledge graph} $\mathcal{K}$ comprises three layers:

\paragraph{Layer 1: Structural Knowledge ($\mathcal{K}_S$).}
The accounting graph topology---accounts as nodes, transaction flows as
directed edges, master data entities with typed relationships, document
chains with referential integrity.  This layer answers \emph{what
entities exist and how they connect}.

\paragraph{Layer 2: Statistical Knowledge ($\mathcal{K}_\Sigma$).}
The distributional properties of all quantities---amount distributions
conforming to Benford's law, temporal patterns (period-end spikes,
intraday profiles, processing lags), cross-field correlations via copula
models, economic cycle dynamics.  This layer answers \emph{what
quantitative patterns characterize normal operation}.

\paragraph{Layer 3: Normative Knowledge ($\mathcal{K}_N$).}
The rules, standards, and controls that govern the enterprise---accounting
framework (US~GAAP, IFRS, etc.), audit standards (ISA, PCAOB), internal
controls (COSO~2013 with 17 principles), regulatory requirements (SOX~302/404),
and organizational policies (approval thresholds, SoD rules).  This layer
answers \emph{what constraints define correct behavior}.
\end{definition}

\begin{proposition}[Reference Knowledge Graph Completeness]
\label{thm:completeness}
Let $\mathcal{G}$ be a fully specified generative model that explicitly
encodes structural rules $\mathcal{R}_S$, statistical parameters
$\Theta_\Sigma$, and normative constraints $\mathcal{C}_N$.  A dataset
$D = \mathcal{G}(\text{seed})$ generated by $\mathcal{G}$ constitutes a
reference knowledge graph that is \emph{complete} with respect to all
three layers:
\begin{equation}
\mathcal{K}(D) = \mathcal{K}_S(\mathcal{R}_S)
  \cup \mathcal{K}_\Sigma(\Theta_\Sigma)
  \cup \mathcal{K}_N(\mathcal{C}_N)
\label{eq:completeness}
\end{equation}
Furthermore, for any predicate $\phi$ over $D$ that is decidable given
$\mathcal{G}$, the ground truth value $\phi(D)$ is computable in time
polynomial in $|D|$.
\end{proposition}

\begin{proof}[Justification]
By construction, every record in $D$ is produced by a deterministic
function of the seed and the generative model $\mathcal{G}$.  The
structural rules $\mathcal{R}_S$ determine which entities exist and
how they relate (Layer~1 completeness).  The statistical parameters
$\Theta_\Sigma$ determine all distributional properties (Layer~2
completeness).  The normative constraints $\mathcal{C}_N$ are explicitly
encoded in the generation process, so every violation or conformance is
known (Layer~3 completeness).  This is a direct consequence of forward
generation from a fully specified model: since the model is known, any
question about the data that depends only on the model's specification
is answerable by inspecting the model.
\end{proof}

% ─────────────────────────────────────────────────────────────────────
\subsection{Forward Generation vs.\ Inverse Recovery}
\label{sec:gt_forward}

We summarize the contrast between the two paradigms:

\begin{center}
\small
\begin{tabular}{p{3.5cm}p{5cm}p{5cm}}
\toprule
\textbf{Property} & \textbf{Inverse Recovery} & \textbf{Forward Generation} \\
\midrule
Ground truth &
  Partially recoverable via domain expertise &
  Known by construction \\
Systematic errors &
  May be inherited if undetected &
  Absent (or injected with labels) \\
Completeness &
  Rich in context but limited by observability &
  Complete across all three layers \\
Complexity &
  Exponential (\Cref{thm:combinatorial}) &
  Linear in output size \\
Verifiability &
  Verification constrained by data availability &
  Every fact is traceable to the model \\
Anomaly labels &
  Typically requires manual annotation &
  Generated with full provenance \\
Reproducibility &
  Depends on data access and agreements &
  Deterministic from seed \\
\bottomrule
\end{tabular}
\end{center}

\noindent The forward generation paradigm does not replace the need to
analyze real enterprise data---it provides the \emph{reference frame}
against which analysis tools, knowledge construction pipelines, and
audit procedures can be developed, tested, and calibrated before
deployment on real data.  \system{} implements this paradigm as a
practical tool, as detailed in the following sections.
